\section*{Capítulo \ref{ch:g5:water-shapes-the-land} --- A água modela a terra}

\begin{solution}{exo:g5:water-shapes-the-land:1}
(a) \omterm{def:g5:water-shapes-the-land:weathering}{Intemperismo}: a \omterm{def:g3:stones-pebbles-sand:rock}{rocha} se quebra no lugar onde está. (b) \omterm{def:g5:water-shapes-the-land:erosion}{Erosão}: a
\omterm{def:g3:stones-pebbles-sand:sand}{areia} é desgastada e levada embora. (c) \omterm{def:g5:water-shapes-the-land:deposition}{Deposição}: a \omterm{def:g3:stones-pebbles-sand:sand}{lama} se deposita. (d)
\omterm{def:g5:water-shapes-the-land:weathering}{Intemperismo}: a raiz quebra a pedra no lugar onde ela está.
\end{solution}

\begin{solution}{exo:g5:water-shapes-the-land:2}
A \omterm{def:g5:water-shapes-the-land:source-mouth}{nascente} é o lugar onde o \omterm{def:g1:landscapes-around-us:water}{rio} começa (uma fonte, uma geleira que
derrete, uma terra alta e encharcada). A \omterm{def:g5:water-shapes-the-land:source-mouth}{foz} é o lugar onde ele termina,
no \omterm{def:g2:water-on-earth:ocean}{mar}, num \omterm{def:g1:landscapes-around-us:water}{lago} ou num \omterm{def:g1:landscapes-around-us:water}{rio} maior. Um \omterm{def:g5:water-shapes-the-land:source-mouth}{afluente} é um \omterm{def:g1:landscapes-around-us:water}{rio} menor que deságua
num \omterm{def:g1:landscapes-around-us:water}{rio} maior.
\end{solution}

\begin{solution}{exo:g5:water-shapes-the-land:3}
A \omterm{def:g5:water-shapes-the-land:source-mouth}{nascente}; o \omterm{def:g1:landscapes-around-us:landform}{vale} estreito em V; os \omterm{def:g5:water-shapes-the-land:meander}{meandros}; o \omterm{def:g5:water-shapes-the-land:delta}{delta}.
\end{solution}

\begin{solution}{exo:g5:water-shapes-the-land:4}
(a) Falso: a geada, a água, as raízes, o calor e o frio quebram devagar
até a \omterm{def:g3:stones-pebbles-sand:rock}{rocha} dura (\omterm{def:g5:water-shapes-the-land:weathering}{intemperismo}). (b) Verdadeiro: quanto mais rápida a
água, maiores os pedaços que ela consegue mover. (c) Falso: os pedaços mais
pesados se depositam primeiro; a \omterm{def:g3:stones-pebbles-sand:sand}{lama} mais fina se deposita por último,
onde a água está quase parada, num \omterm{def:g1:landscapes-around-us:water}{lago} ou no \omterm{def:g2:water-on-earth:ocean}{mar}.
\end{solution}

\begin{solution}{exo:g5:water-shapes-the-land:5}
A fica do lado de fora da curva: ali a água é rápida e funda, e ela
desgasta a margem (\omterm{def:g5:water-shapes-the-land:erosion}{erosão}). B fica do lado de dentro da curva: ali a água
é lenta, e ela deixa \omterm{def:g3:stones-pebbles-sand:sand}{areia} (\omterm{def:g5:water-shapes-the-land:deposition}{deposição}).
\end{solution}

\begin{solution}{exo:g5:water-shapes-the-land:6}
Nas \omterm{def:g1:landscapes-around-us:landform}{montanhas}, o \omterm{def:g1:landscapes-around-us:water}{rio} é íngreme e rápido: ele consegue rolar pedras grandes
e leva os grãos pequenos embora. Perto do \omterm{def:g2:water-on-earth:ocean}{mar}, o \omterm{def:g1:landscapes-around-us:water}{rio} é lento: ele não
consegue mais mover pedras grandes, só \omterm{def:g3:stones-pebbles-sand:sand}{areia} e \omterm{def:g3:stones-pebbles-sand:sand}{lama}, e as deixa ali. No
caminho, além disso, as pedras foram quebradas e desgastadas em grãos
cada vez menores.
\end{solution}

\begin{solution}{exo:g5:water-shapes-the-land:7}
(a) Curso superior: $600 \div 6 = 100$ quilômetros. Curso médio:
$600 \div 2 = 300$ quilômetros. (b) Curso inferior: $600 - 100 - 300 = 200$
quilômetros. Como $3 \times 200 = 600$, ele é $\frac{1}{3}$ do \omterm{def:g1:landscapes-around-us:water}{rio}.
(Verificação: $\frac{1}{6} + \frac{1}{2} + \frac{1}{3} = \frac{1}{6} +
\frac{3}{6} + \frac{2}{6} = \frac{6}{6}$.)
\end{solution}

\begin{solution}{exo:g5:water-shapes-the-land:8}
As manchas claras são \omterm{def:g3:stones-pebbles-sand:sand}{areia} e \omterm{def:g3:stones-pebbles-sand:pebble}{cascalho} depositados do lado de dentro das
curvas, onde a água é lenta (\omterm{def:g5:water-shapes-the-land:deposition}{deposição}). Os \omterm{def:g1:landscapes-around-us:water}{lagos} curvos e escuros são
curvas antigas do \omterm{def:g1:landscapes-around-us:water}{rio}: duas curvas cresceram até ficar muito próximas, o
\omterm{def:g1:landscapes-around-us:water}{rio} rompeu a faixa de terra entre elas, pegou o atalho e deixou a curva
antiga para trás, como um \omterm{def:g1:landscapes-around-us:water}{lago}.
\end{solution}

\begin{solution}{exo:g5:water-shapes-the-land:9}
Os pastos planos ao lado do \omterm{def:g1:landscapes-around-us:water}{rio} são a sua \omterm{def:g5:water-shapes-the-land:flood}{planície de inundação}: a terra
que o \omterm{def:g1:landscapes-around-us:water}{rio} cobre quando transborda. Eles são planos justamente porque o
\omterm{def:g1:landscapes-around-us:water}{rio} os construiu com a sua própria \omterm{def:g3:stones-pebbles-sand:sand}{lama}. Um dia o \omterm{def:g1:landscapes-around-us:water}{rio} vai transbordar de
novo e encher a escola de água, e a água rápida de uma enchente é
perigosa. A escola deve ser construída num terreno mais alto, e os pastos
devem continuar sendo campos, pastos ou um parque.
\end{solution}

\begin{solution}{exo:g5:water-shapes-the-land:10}
(a) Um ano tem 12 meses: $12$ bilhões $\div\, 12 = 1$ bilhão de toneladas
por mês, em média. (b) Um bilhão é $1\,000\,000\,000$; dividido por 10, dá
$100\,000\,000$: cem milhões de cargas de caminhão.
\end{solution}

\begin{solution}{exo:g5:water-shapes-the-land:11}
Uma história possível. O grão faz parte de uma \omterm{def:g3:stones-pebbles-sand:rock}{rocha} no alto de uma
\omterm{def:g1:landscapes-around-us:landform}{montanha}. A geada abre uma rachadura em volta dele, noite após noite, até
que um pedaço de \omterm{def:g3:stones-pebbles-sand:rock}{rocha} se solta: isso é o \omterm{def:g5:water-shapes-the-land:weathering}{intemperismo}. O pedaço cai no
riacho lá embaixo. No curso superior, a água é rápida e íngreme: ela faz o
pedaço rolar (\omterm{def:g5:water-shapes-the-land:erosion}{erosão}) e, batendo em outras pedras, ele se quebra em \omterm{def:g3:stones-pebbles-sand:pebble}{seixos}
e depois em grãos de \omterm{def:g3:stones-pebbles-sand:sand}{areia}. No curso médio, o \omterm{def:g1:landscapes-around-us:water}{rio} fica mais lento; o grão
descansa por um tempo na praia de \omterm{def:g3:stones-pebbles-sand:sand}{areia} do lado de dentro da curva de um
\omterm{def:g5:water-shapes-the-land:meander}{meandro}, até que uma cheia o pega de novo. No curso inferior, o \omterm{def:g1:landscapes-around-us:water}{rio} é
lento e largo; o grão viaja com a água até a \omterm{def:g5:water-shapes-the-land:source-mouth}{foz}. Ali a água quase para,
e o grão se deposita com a \omterm{def:g3:stones-pebbles-sand:sand}{lama} e a \omterm{def:g3:stones-pebbles-sand:sand}{areia} do \omterm{def:g5:water-shapes-the-land:delta}{delta}: isso é a \omterm{def:g5:water-shapes-the-land:deposition}{deposição}.
\end{solution}
